\chapter{Introduction}
\label{chap:intro}
\chaptermark{Optional running chapter heading}

\section{Background}
\label{intro:bgr}

Interactive Theorem Provers\footnote{Commonly abbreviated as ITPs} are domain-specific languages
designed for formalizing mathematics.
They are essential for modern mathematics research,
granting a near-absolute degree of certainty by verifying proofs via deterministic algorithms and
managing extreme complexity by decomposing large-scale proofs into reusable components.
Therefore, they are frequently utilized by mathematicians
who may lack a comprohensive programming background to employ ITPs effectively.

Integrated Development Environment\footnote{Commonly abbreviated as IDEs} integration with ITP
represents a viable approach to addressing this challenge by
equipping user with rapid, real-time code diagnostics, seamless code navigation,
syntax highlighting, and other advanced features.

\section{Research Gap}
\label{intro:gap}

Rzk is a theorem prover based on Riehl and Shulman's type theory\cite{riehl2023typetheorysyntheticinftycategories} with shapes.
While its existing language server keeps track of unsaved changes,
highlights syntax and enables navigation, it suffers from several limitations:
proof evaluation is not automatically executed upon file modification,
and large projects exhibit significant latency during evaluation
when changes are introduced at the beginning of the codebase.

\section{Approach}
\label{intro:plan}

To achieve these objectives, the research presented in this thesis is structured in several consecutive stages.

Initially, the typechecker module of language server will be adapted
to leverage unsaved changes from the IDE where possible, thereby providing relevant diagnostics.

Subsequently, an edit-triggered evaluation policy will be introduced
to deliver these diagnostics automatically without requiring explicit manual triggers.

These enhancements establish the foundation for the core contribution of this thesis,
enabling the evaluation of diagnostic delivery latencies
within the current architecture of the Rzk LSP\footnote{Language Server Protocol} server.

Following this, a build system incorporating module dependencies and cache invalidation
will be designed and implemented to eliminate redundant type checking of unmodified code,
thereby minimizing the latency of diagnostic delivery.


\section{Contributions}
\label{intro:contr}

In summary, this thesis delivers several contributions that optimize IDE behavior as follows:
\begin{itemize}
	\item The language server provides relevant diagnostics for unsaved changes directly within the IDE.
	In this context, "relevant diagnostics" implies that diagnostics for unsaved changes matches
	the output of the Rzk typechecker executed on the corresponding source files on disk.
	\item When user begins typing, outdated diagnostic messages are immediately suppresed.
	Following a 500 ms timeout after the final keystroke, the language server triggers typecheck on the project,
	presenting results as fresh dianostics.
	This debouncing strategy allows language server to initiate computations promptly
	while mitigating the overhead of repeatedly halting and restarting the checker during rapid continuous typing.
	\item The language server minimizes computational overhead by avoiding repetitive execution on unchanged code.
	When a user modifies a module, typecheckind is isolated to that specific module and
	its downstream dependencies, rather than every subsequent file in the project hierarchy.
	Although this optimization is less noticeable to the user, this strategy is expected
	to substantially reduce re-evaluation latencies for large projects when early-stage files are modified,
	while maintaining standard performance for changes made in later files.
	For instance, the evaluated performance metrics of Rzk and its language server on
	the simplicial HoTT formalization project (sHoTT)\cite{shott} -
	which comprises 13 files with rzk code - are present in Table \ref{tab:perf}.
\end{itemize}

\begin{table}[b]
	\centering
	\begin{center}
		\begin{tabular}{|c|c|c|c|}
			\hline
			& Whole project & 00-common.rzk.md & 13-truncation-levels \\ \hline
			Rzk & 3000 & 200 & 300 \\ \hline
			Current language server & 3000 & 3000 & 300 \\ \hline
			Improved language server & 3000 & 200 & 300 \\ \hline
		\end{tabular}
	\end{center}
	\caption{Evaluted diagnostic latency of Rzk, its current language server and its implementation proposed in this thesis on the sHoTT project (computation time in ms)}
	\label{tab:perf}
\end{table}

\section{Paper Outline}
\label{intro:outline}

The remainder of this paper is structured as follows.
Chapter 2 reviews Rzk, the LSP, and comparable solutions implemented in other ITPs.
Chapter 3 describes the applied methodology.
Chapter 4 documents the design and implementation aspects.
Chapter 5 presents and analyzes the experimental results.
Chapter 6 concludes the paper and summarizes the completed work.
